\begin{codebox}
\codeline{\textcolor{cmtgray}{\#\ 单调与非单调推理示例}}
\codeline{\textcolor{cmtgray}{\#\ Monotonic\ and\ Non-monotonic\ Reasoning\ Examples}}
\codeline{\textcolor{cmtgray}{\#}}
\codeline{\textcolor{cmtgray}{\#\ 推理是本体的核心能力：从显性知识推导隐性知识}}
\codeline{\textcolor{cmtgray}{\#\ 单调性（Monotonicity）是划分推理机制的重要维度}}
\codeblank{}
\codeline{\textcolor{cmtgray}{\#\ ==============================}}
\codeline{\textcolor{cmtgray}{\#\ 1.\ 单调推理\ (Monotonic\ Reasoning)}}
\codeline{\textcolor{cmtgray}{\#\ ==============================}}
\codeline{\textcolor{cmtgray}{\#\ 定义：新增知识只会增加结论，不会推翻已有结论}}
\codeline{\textcolor{cmtgray}{\#\ 经典逻辑（一阶逻辑、描述逻辑）都是单调的}}
\codeline{\textcolor{cmtgray}{\#\ OWL\ DL\ 推理器（Pellet\ /\ HermiT）遵循单调性假设}}
\codeblank{}
\codeline{\textcolor{cmtgray}{\#\ 已知公理（TBox）：}}
\codeline{CNCLathe\ \(\sqsubseteq\)\ CNCMachine\ \ \ \ \ \ \ \ \ \ \ \ \ \#\ 数控车床是数控机床}
\codeline{CNCMachine\ \(\sqsubseteq\)\ ProcessingEquipment\ \ \#\ 数控机床是加工设备}
\codeblank{}
\codeline{\textcolor{cmtgray}{\#\ 已知事实（ABox）：}}
\codeline{CNCLathe(Lathe\_001)\ \ \ \ \ \ \ \ \ \ \ \ \ \ \ \#\ Lathe\_001\ 是一台数控车床}
\codeblank{}
\codeline{\textcolor{cmtgray}{\#\ 单调推导（子类关系传递）：}}
\codeline{CNCMachine(Lathe\_001)\ \ \ \ \ \ \ \ \ \ \ \ \ \#\ 推论1：Lathe\_001\ 是数控机床}
\codeline{ProcessingEquipment(Lathe\_001)\ \ \ \ \#\ 推论2：Lathe\_001\ 是加工设备}
\codeblank{}
\codeline{\textcolor{cmtgray}{\#\ 单调性保证：之后无论向知识库添加什么新事实}}
\codeline{\textcolor{cmtgray}{\#\ （设备故障、位置变更、新的分类公理），}}
\codeline{\textcolor{cmtgray}{\#\ 推论1、推论2\ 永远成立，不会被撤销}}
\codeblank{}
\codeline{\textcolor{cmtgray}{\#\ ==============================}}
\codeline{\textcolor{cmtgray}{\#\ 2.\ 非单调推理\ (Non-monotonic\ Reasoning)}}
\codeline{\textcolor{cmtgray}{\#\ ==============================}}
\codeline{\textcolor{cmtgray}{\#\ 定义：新增知识可能推翻已有结论（可撤销推理，Defeasible\ Reasoning）}}
\codeline{\textcolor{cmtgray}{\#\ 适用于"默认假设\ +\ 例外"形式的工程知识}}
\codeblank{}
\codeline{\textcolor{cmtgray}{\#\ 默认规则（Reiter\ 默认逻辑）：}}
\codeline{Equipment(x)\ :\ available(x)\ /\ available(x)}
\codeline{\textcolor{cmtgray}{\#\ 规则解读：如果\ x\ 是设备，且"x\ 可用"与知识库一致（无矛盾信息），}}
\codeline{\textcolor{cmtgray}{\#\ 则默认推出\ x\ 可用}}
\codeblank{}
\codeline{\textcolor{cmtgray}{\#\ 阶段1\ ——\ 知识库仅有：}}
\codeline{Equipment(Lathe\_001)}
\codeline{\textcolor{cmtgray}{\#\ 默认推论：available(Lathe\_001)\ 成立\ \(\checkmark\)}}
\codeblank{}
\codeline{\textcolor{cmtgray}{\#\ 阶段2\ ——\ 新增事实与规则：}}
\codeline{hasFault(Lathe\_001)\ \ \ \ \ \ \ \ \ \ \ \ \ \ \ \#\ Lathe\_001\ 出现故障}
\codeline{\(\forall\)x:\ hasFault(x)\ \(\rightarrow\)\ \(\lnot\)available(x)\ \ \ \#\ 故障设备不可用}
\codeline{\textcolor{cmtgray}{\#\ 撤销推论：available(Lathe\_001)\ 被推翻\ \(\times\)}}
\codeline{\textcolor{cmtgray}{\#\ 对于没有故障记录的其他设备，仍默认其可用}}
\codeline{\textcolor{cmtgray}{\#\ 这种"结论收回"是单调逻辑无法表达的}}
\codeblank{}
\codeline{\textcolor{cmtgray}{\#\ ==============================}}
\codeline{\textcolor{cmtgray}{\#\ 3.\ 开放世界假设与封闭世界假设}}
\codeline{\textcolor{cmtgray}{\#\ ==============================}}
\codeline{\textcolor{cmtgray}{\#\ OWA\ (Open\ World\ Assumption)\ \ ——\ OWL\ /\ 描述逻辑采用}}
\codeline{\textcolor{cmtgray}{\#\ \ \ 未声明的事实\ =\ 未知（可能为真，也可能为假）}}
\codeline{\textcolor{cmtgray}{\#\ CWA\ (Closed\ World\ Assumption)\ ——\ 关系数据库\ /\ Prolog\ 采用}}
\codeline{\textcolor{cmtgray}{\#\ \ \ 未声明的事实\ =\ 假（否定即失败，Negation\ as\ Failure）}}
\codeblank{}
\codeline{\textcolor{cmtgray}{\#\ 示例：知识库中只有一条事实}}
\codeline{producedBy(Product\_A,\ Lathe\_001)}
\codeblank{}
\codeline{\textcolor{cmtgray}{\#\ 问题：Product\_A\ 是否由\ Mill\_002\ 生产？}}
\codeline{\textcolor{cmtgray}{\#\ OWA\ 回答：未知（除非显式声明\ \(\lnot\)producedBy(Product\_A,\ Mill\_002)）}}
\codeline{\textcolor{cmtgray}{\#\ CWA\ 回答：否（知识库未记录即视为假）}}
\codeblank{}
\codeline{\textcolor{cmtgray}{\#\ 工程影响：}}
\codeline{\textcolor{cmtgray}{\#\ 在\ OWL\ 中校验"设备未分配任何任务"这类否定条件时，}}
\codeline{\textcolor{cmtgray}{\#\ 需借助\ SPARQL\ 的\ FILTER\ NOT\ EXISTS（局部封闭世界）实现}}
\codeblank{}
\codeline{\textcolor{cmtgray}{\#\ ==============================}}
\codeline{\textcolor{cmtgray}{\#\ 4.\ 工程应用对比}}
\codeline{\textcolor{cmtgray}{\#\ ==============================}}
\codeline{\textcolor{cmtgray}{\#\ 单调推理适用场景：}}
\codeline{\textcolor{cmtgray}{\#\ \ \ 类层次分类、一致性检查、设备能力匹配}}
\codeline{\textcolor{cmtgray}{\#\ \ \ （要求结论稳定可靠、可审计）}}
\codeline{\textcolor{cmtgray}{\#\ 非单调推理适用场景：}}
\codeline{\textcolor{cmtgray}{\#\ \ \ 默认配置（设备默认可用）、异常处理（故障撤销可用性）、}}
\codeline{\textcolor{cmtgray}{\#\ \ \ 计划修订（新约束到达后撤销原有调度结论）}}
\codeblank{}
\codeline{\textcolor{cmtgray}{\#\ 实现说明：OWL\ 标准推理器只支持单调推理；}}
\codeline{\textcolor{cmtgray}{\#\ 非单调能力需在应用层借助规则引擎实现，}}
\codeline{\textcolor{cmtgray}{\#\ 如\ SPARQL\ 更新、SHACL\ 规则、Jena\ Rules，}}
\codeline{\textcolor{cmtgray}{\#\ 或支持否定即失败（NAF）的\ Datalog\ 引擎}}
\end{codebox}
